% Part: incompleteness
% Chapter: arithmetization-syntax
% Section: introduction

\documentclass[../../../include/open-logic-section]{subfiles}

\begin{document}

\olfileid{inc}{art}{int}
\olsection{Pambuka}

Kanggo ngubungake komputabilitas lan logika,
kita mbutuhake cara kanggo ngrembug objek logika
(simbol, suku, !!{formula}s, !!{derivation}s),
operasi marang objek-objek iku, sarta sifat
lan relasine kanthi cara kang bisa diolah
sacara komputasional. Iki bisa ditindakake
langsung kanthi nggatekake fungsi lan relasi
komputabel ing simbol, urutan simbol, lan
objek liyane kang diwangun saka simbol.
Amarga kabeh objek sintaks logis winates
lan diwangun saka himpunan simbol kang
!!a{enumerable}, cara iki bisa digunakake
kanggo sawetara model komputasi. Nanging
model komputasi liyane---kayata fungsi
rekursif---winates mung ing wilangan,
relasi, lan fungsine. Kajaba iku, pungkasane
kita uga kepengin ngolah sintaks ing njero
teori tartamtu, mligine teori kang dirumusake
ing basa aritmetika. Ing kahanan iki, kita
kudu \emph{ngaritmetisasi} sintaks: nggambarake
objek sintaks, operasi marang objek kasebut,
lan relasine minangka wilangan, fungsi
aritmetika, lan relasi aritmetika,
miturut urutane. Gagasan kang wiwit saka
Leibniz iki yaiku menehi angka marang
objek sintaks.

Menehi wilangan minangka ``kode'' kanggo
simbol iku cukup gampang. Sawetara simbol
rada njlimet, amarga, umpamane, ana tanpa
wates akeh !!{variable}s, lan malah tanpa
wates akeh !!{function}s kanggo saben
aritas~$n$. Nanging mesthi bisa menehi
wilangan marang simbol kanthi sistematis,
saengga, umpamane, $\Obj v_2$ lan
$\Obj v_3$ nduwé kode kang béda. Urutan
simbol (kayata suku lan !!{formula}s)
luwih njlimet maneh. Nanging yen urutan
wilangan bisa diolah nganggo aritmetika
wae (umpamane nganggo pangodean urutan
liwat pangkat wilangan prima), pangodean
simbol siji-siji bisa diperluas dadi
pangodean urutan simbol, banjur dadi
pangodean urutan utawa susunan liyane
saka !!{formula}s, kayata !!{derivation}s.
Pangodean kang diperluas iki diarani
``panomeran G\"odel''. Saben suku,
!!{formula}, lan !!{derivation} diwenehi
wilangan G\"odel.

Kanthi ngodeake urutan simbol minangka
urutan kode-kodene, lan milih sistem
pangodean urutan kang bisa diolah nganggo
fungsi komputabel, kita banjur bisa ngolah
wilangan G\"odel nganggo fungsi komputabel.
Ing praktik, kabeh fungsi kang dibutuhake
bakal rekursif primitif. Tuladhane,
ngitung dawa urutan lan ngitung unsur
kapangka-$i$ saka sawijining urutan adhedhasar
kode urutane iku rekursif primitif.
Yen wilangan kang ngodeake urutan iku,
umpamane, wilangan G\"odel saka
!!a{formula}~$!A$, langsung katon yen
dawa !!a{formula} lan (kode saka) simbol
kapangka-$i$ ing !!a{formula} uga bisa
diitung saka wilangan G\"odel~$!A$.
Luwih angel sethithik mbuktekake yen,
umpamane, sifat minangka wilangan G\"odel
saka suku kang kabentuk kanthi bener utawa
saka !!{derivation} kang bener iku rekursif
primitif. Nanging iki bisa ditindakake,
amarga urutan kang dirembug (suku,
!!{formula}s, !!{derivation}s)
ditegesi kanthi induktif.

Minangka tuladha, gatekna operasi substitusi.
Yen $!A$ iku formula, $x$ variabel, lan
$t$ suku, mula $\Subst{!A}{t}{x}$ iku
asil ngganti saben panggonan bebas~$x$
ing~$!A$ nganggo~$t$. Saiki upamane kita
wis menehi wilangan G\"odel marang $!A$,
$x$, $t$---umpamane $k$, $l$, lan $m$
miturut urutane. Skema kang padha menehi
wilangan G\"odel marang
$\Subst{!A}{t}{x}$, umpamane~$n$.
Peta saka $k$, $l$, lan $m$ menyang
$n$ iki iku analog aritmetika saka operasi
substitusi. Nalika operasi substitusi
ngirim $!A$, $x$, $t$ menyang
$\Subst{!A}{t}{x}$, fungsi substitusi
kang diaritmetisasi ngirim wilangan G\"odel
$k$, $l$, $m$ menyang wilangan G\"odel~$n$.
Kita bakal ndeleng yen fungsi iki rekursif
primitif.

Aritmetisasi sintaks ora mung narik kawigaten
minangka gagasan abstrak. Ing wiwitan, iki
pancen sawijining panemu kang ora prasaja:
basa kaya basa aritmetika, kang ora nduwé
piranti kanggo ``ngrembug'' basa, jebul
bisa ngformalake sifat-sifat rumit saka
ekspresi. Saka kene, tinggal takon apa
kang bisa \emph{dibuktekake} dening teori
ing basa iki, kayata aritmetika Peano,
ngenani basane dhewe (kalebu, umpamane,
apa !!{sentence}s bisa dibuktekake utawa
bener). Iki nggawa kita menyang teorema
watesan kang misuwur saka G\"odel
(ngenani ora bisa dibuktekake) lan Tarski
(ngenani kabeneran kang ora bisa ditegesi).
Nanging cara ngaritmetisasi sintaks uga
wigati kanggo mbuktekake sawetara asil
penting ing teori komputabilitas, umpamane
ngenani kakuwatan komputasional teori
utawa sesambungan antarane model-model
komputabilitas kang béda. Aritmetisasi
sintaks dadi tuladha kanggo ngaritmetisasi
objek lan sifat liyane. Tuladhane,
konfigurasi lan komputasi (umpamane
saka mesin Turing) uga bisa diaritmetisasi.
Iki ndadekake komputasi ing sawijining
model (umpamane mesin Turing) bisa
disimulasi ing model liyane (umpamane
fungsi rekursif).

\end{document}
